\section*{Chapter \ref{ch:fm:culture-and-decision-making} --- Culture and Decision-Making}

\begin{solution}{exo:fm:culture-and-decision-making:1}
$(0.09+0.09+0.49)/3=0.223$; a constant 0.5 scores 0.25.
\end{solution}

\begin{solution}{exo:fm:culture-and-decision-making:2}
0.70 for the probabilities; the log-odds average gives 0.710, further from one half.
\end{solution}

\begin{solution}{exo:fm:culture-and-decision-making:3}
By 5 when independent; with $\rho=0.5$ the variance is $0.5+0.5/5=0.6$ of one forecaster's, a factor of only 1.67.
\end{solution}

\begin{solution}{exo:fm:culture-and-decision-making:4}
Its reliability term is 0.040 against 0.017: it stretches its log-odds, so its extreme forecasts come true far less often than stated; its resolution
is about the same.
\end{solution}

\begin{solution}{exo:fm:culture-and-decision-making:5}
The forecasters are already overconfident, so pushing the average further from one half adds miscalibration. Extremising helps when forecasters
hold different information and are individually well calibrated, so that their average is too moderate.
\end{solution}

\begin{solution}{exo:fm:culture-and-decision-making:6}
Collect each person's probability that the strategy is broken before the meeting, discuss the evidence, collect again, and record the aggregate
of the log-odds in the journal with the kill criterion of chapter 8.
\end{solution}

\begin{solution}{exo:fm:culture-and-decision-making:7}
Of 362 bets with positive expected value by the forecaster's own probabilities, 123 lost: an outcome-based review would call 34\% of them mistakes.
\end{solution}

\begin{solution}{exo:fm:culture-and-decision-making:8}
A winning outcome says little about the decision: judge it by the probabilities and information at the time, recorded in the journal, and by the
calibration of the person's past forecasts.
\end{solution}

\begin{solution}{pb:fm:culture-and-decision-making:1}
\begin{enumerate}
\item Probability training, teaming and tracking improved calibration and resolution; with aggregation, the best forecasts two years in a row.
\item See \cref{def:fm:culture-and-decision-making:journal}.
\item See \cref{def:fm:culture-and-decision-making:brier}; 0.25.
\item Reliability minus resolution plus uncertainty: miscalibration, discrimination between events, and the events' own unpredictability.
\item 0.240, 0.225, 0.252, 0.245 and 0.262; the truth 0.214.
\item Forecaster 1: reliability 0.017, resolution 0.027; forecaster 5: 0.040 and 0.026; the log-odds average: 0.010 and 0.039, with uncertainty 0.25.
\item Probability mean 0.220, log-odds mean 0.223, extremised 0.237.
\item The overconfident forecaster's curve is flatter than the diagonal; the average lies close to it.
\item See \cref{def:fm:culture-and-decision-making:agg}.
\item See \cref{prop:fm:culture-and-decision-making:rho}.
\item 0.207 independent; 0.224 at 0.5; 0.237 at 1.
\item People anchor on the first number, the loudest argument or the senior person's view.
\item Independent numbers first, discussion second, a second round, aggregation of the log-odds.
\item See \cref{def:fm:culture-and-decision-making:outcome}; 34\% of positive-expected-value bets lost, and forecaster 2 won 66\% while forecasting
77.7\%.
\item See \cref{def:fm:culture-and-decision-making:premortem}.
\item Through scores over long periods and many forecasts, not single outcomes (chapter 10).
\item The breach, the options, the probability that the position recovers, the expected cost of each option, who approved, and the review date.
\item Gaming by forecasting only easy events, hedging towards one half, and discouraging dissent; score all significant decisions, with
calibration and resolution.
\item In the chapter's model, aggregating five independent forecasts scores 0.207 against 0.224 for one discussed forecast at $\rho=0.5$ (0.017
better) and 0.237 at $\rho=1$ (0.030 better).
\item Collect written probabilities from each person before any discussion, aggregate their log-odds and score everyone quarterly; keep the
meeting, but after the numbers are in.
\end{enumerate}
\end{solution}

\begin{solution}{iq:fm:culture-and-decision-making:1}
The mean squared error of probability forecasts of binary outcomes; 0.25 is what always saying one half scores.
\iqlookfor{the no-skill benchmark.}
\end{solution}

\begin{solution}{iq:fm:culture-and-decision-making:2}
Look at the probabilities and information recorded when you decided and at your calibration over many such trades, not at this outcome.
\iqlookfor{process over outcome.}
\end{solution}

\begin{solution}{iq:fm:culture-and-decision-making:3}
Log-odds add evidence symmetrically and are not dragged towards one half by a single moderate forecaster; averaging probabilities is too
conservative when forecasters hold different information.
\iqlookfor{evidence adds in log-odds.}
\end{solution}

\begin{solution}{iq:fm:culture-and-decision-making:4}
To $0.5+0.5/5=0.6$ of one person's: a 40\% reduction, against 80\% if independent.
\iqlookfor{$\rho+(1-\rho)/K$.}
\end{solution}

\begin{solution}{iq:fm:culture-and-decision-making:5}
Before committing, have the team assume the entry failed after a year and write down why; rank the reasons, add the likely ones to the plan's
kill criteria and critical path.
\iqlookfor{imagined failure, then action.}
\end{solution}

\begin{solution}{iq:fm:culture-and-decision-making:6}
Group forecasts into bins; expand $\frac1N\sum(p_i-y_i)^2$ around each bin's observed frequency and around the overall frequency; with forecasts
constant within bins the cross terms vanish, leaving reliability minus resolution plus uncertainty.
\iqlookfor{the binning and the vanishing cross terms.}
\end{solution}
